\begin{tcolorbox}[
    colback=SkyBlue!5,
    colframe=blue!20,
    boxrule=1pt,
    arc=2pt,
    width=\textwidth,
    boxsep=5pt,
    left=5pt,
    right=5pt,
    top=5pt,
    bottom=5pt
]

\noindent
\begin{minipage}[t]{0.7\textwidth}
\begin{example}
Consider the following program. It has 4 execution graphs: 
one where the line $a_2:=x$ reads from the $\init$, one where it reads from $x:=1$, and two where it reads from $x:=2$. 
The algorithm in \cite{gator} runs as follows. 
It first schedules $t_1$ and writes 1 to $x$. 
Then since both threads are enabled with a read, the scheduler picks one uniformly at random.
A case analysis of the two possibilities shows that the expected value of the estimator is
3.5.
\end{example}
\end{minipage}%
\hfill
\begin{minipage}[t]{0.3\textwidth}
\vspace{0pt}
\begin{equation*}
  \tag{\textsc{wrww+rr}}
  \label{ex:wrww+rr}
  \inarr{ 
\writeInst{x}{1} \\
\readInst{a_1}{x} \\
\writeInst{x}{2} \\
\writeInst{y}{1} 
}{ 
\readInst{a_2}{x} \\
\text{if } a_2 = 2 \\
\readInst{a_2}{y }}
\end{equation*}
\end{minipage}%

\begin{comment}
\vspace{0.3cm}

Case A: the scheduler proceeds with $t_1$. Since the remaining events of $t_1$ are writes, the scheduler will continue proceeding with $t_1$ until it finishes. Up to this point, the estimation is $1$ because no forward revisit is needed. It then schedules the first event of $t_2$. There are three possible read-from values (including the initial value), so the estimator will store the knowledge that the read $a_2:=x$ has three possible forward revisits. With probability $2/3$, the estimator explores the case that $a_2$ does not read $2$, then $t_2$ terminates and the estimator returns 3. With probability $1/3$, $a_2$ reads  $2$ and enters the if-branch and reads $y$, which has two possible values, so the estimator returns $6$.

Case B: the scheduler proceeds with $t_2$. At that moment, the read has two possible forward revisits (from $\init$ and from the write $x:=1$), so the estimator records $2$ for this read. Regardless of which value $a_2$ reads in this case, it does not enter the if-branch and terminates. $t_1$ is then scheduled, and the estimator detects a backward revisit and updates that the earlier read into $a_2$ actually had three possible values to read from; however, because the algorithm forbids backward revisits, it does not explore that alternative and still returns $3$.

Thus the expected output in Case A is $(2/3)\times3 + (1/3)\times6 = 4$, and in Case B it is $3$. 
Thus, the expected value of the estimator is $0.5\times 4 + 0.5\times 3 = 3.5$.
\end{comment}
\qed
\end{tcolorbox}
